\chapter{How to give credit}
\section{Wish list}
\begin{itemize}
\item Send me a reprint if you used {\tt Migrate} for 
your publication [people obviously do not read that far in this
documentation, because I received only about 15 until July 2010]
\item Cite the documentation and our papers, see below.
\item Report problems to {\tt beerli@fsu.edu}
\item Suggestions (if you need these improvements very soon,  
add a check so that I can hire a programmer to implement 
all those \includegraphics[height=12pt]{mim/smiley})
\end{itemize}

\section{How to give credit}
Maximum likelihood: Beerli 1998; Beerli and Felsenstein 1999, 2001.\\
Bayesian inference: Beerli 2006; Beerli and Felsenstein 2001. \\
Missing population issues: Beerli 2004.\\
General use of \migrate: Beerli 2009.\\
Bayes factor and marginal likelihood: Beerli and Palczewski 2010.\\

\begin{description}
%\bibitem{beerli2004-827}
\lit{Beerli, P.} (1998)
Estimation of migration rates and population sizes in geographically
structured populations. 
In: Advances in molecular ecology (Ed. G. Carvalho). NATO-ASI workshop series. IOS Press, Amsterdam. Pp. 39-53. 

\lit{Beerli, P. and J.~Felsenstein} (1999)
Maximum-likelihood estimation of migration rates and effective
  population numbers in two populations using a coalescent approach.
{\em Genetics}, 152(2):763--73, 1999


\lit{Beerli, P. and J.~Felsenstein} (2001)
Maximum likelihood estimation of a migration matrix and effective
  population sizes in n subpopulations by using a coalescent approach.
{\em Proceedings of the National Academy of Sciences of the USA},
  98(8):4563--4568.

\lit {Beerli, P.} (2004)
Effect of unsampled populations on the estimation of population sizes
  and migration rates between sampled populations.
 {\em Molecular Ecology}, 13:827--836.

\lit{Beerli, P.} (2006)
Comparison of {B}ayesian and maximum likelihood inference of
  population genetic parameters.
{\em Bioinformatics}, 22(3):341--345.

\lit{Beerli, P.} (2009)
How to use migrate or why are markov chain monte carlo programs
  difficult to use? In G.~Bertorelle, M.~W. Bruford, H.~C. Hauffe, A.~Rizzoli, and
  C.~Vernesi, editors, {\em Population Genetics for Animal Conservation},
  volume~17 of {\em Conservation Biology}, pages 42--79. Cambridge University
  Press, Cambridge UK, 2009.

\lit{Beerli, P. and M.~Palczewski}(2010)
Unified framework to evaluate panmixia and migration direction among multiple sampling locations.
{\em Genetics}, 185:313--326

\end{description}


\section{Copyright}
(c) Copyright 1996-2003 by Peter Beerli and Joseph Felsenstein, Seattle WA, USA. 
(c) Copyright 2004-2010 by Peter Beerli, Tallahassee, FL, USA
Permission is granted to copy this document and the program {\it Migrate-n}
and {\it Migrate} 
provided that no fee is charged for it and that this copyright notice is not removed.                               
 
\section{Acknowledgement}
This project is and was supported by grants from National Science 
Foundation (USA) BIR 9527687 and National Health Institutes (USA)
GM51929 and HG01989 all to Joseph Felsenstein and a fellowship of the Swiss National 
Science Foundation to Peter Beerli (1994-1996). 
Development of parts of the MPI structures was funded through a First-Year-Summer-Grant at Florida State University 2004.
Bayes inference development, Bayes factors, data type additions are supported by the joint NSF/NIGMS Mathematical Biology 
program under NIH grant R01 GM 078985.
\par
I thank Mary K. Kuhner, Jon Yamato, Michal Palczewski, and Koffi Sampson for help during debugging and many discussion.
\par
I thank all the people who thought it  worth to report errors and 
foggyness in menu and explanation.

\newpage
